\cotefile{11}
\scanpage{29}
% L1–2 | top
Donc on trouve un foncteur \deleted{\unsure{contravariant}}
\[
\unsure{\mathcal V|V} \longrightarrow \text{Algèbres bivariantes jaugées}
\]

% L3–6 | upper
3) Dans $\mathcal V^{\unsure{o|0|\circ}}$ \deleted{\unsure{il y a un}} il y a objet initial : $k$.

$\unsure{\forall|\mathrm{tt}|\text{Si}}\, A,B\in\mathcal V^{o}$, il existe un coproduit dans $\mathcal V^{o}$ (\deleted{produit} ds $\mathcal V$) savoir le produit tensoriel gauche $A\otimes B$. Soit \deleted{\unsure{les}} $f^*:A\unsure{-|\to}A'$, $g^*:B-B'$, alors $f^*\otimes g^*:A\otimes B\to A'\otimes B'$

% L7 | upper middle
défini par
\[
f\otimes g^*(a\otimes b)=f^*(a)\otimes g^*(b)
\]

% L8–12 | upper middle, additions on both sides of the display
et il \unsure{leur|hom.} correspond\unsure{ent|ant} des coproduits. \deleted{\unsure{On les note}} On convient de mettre sur $A\otimes B$ l'augmentation définie par
\[
k_{A\otimes B}=k_A\otimes k_B
\]
(produit tensoriel \uline{gauche} d'homs homogènes d'algèbres graduées)
\margin{i.e. $k_{A\otimes B}(x\otimes y)=k_A(x)k_B(\unsure{x|y})(-1)^{Bx}$}
\deleted{\unsure{je dis que}} l'on \deleted{\unsure{ait}} alors \deleted{\unsure{On aurait}}
\[
\varphi(f^*\otimes g^*)=\varphi(f^*)\otimes\varphi(g^*)
\]
\margin{produit tensoriel \uline{gauche} d'homs. \uline{homogènes} des \deleted{\unsure{groupes|produits}} modules gradués}

% L13–16 | centre
i.e
\[
(f\times g)_*(a'\otimes b')=f_*(a')\otimes g_*(b')\,(-1)^{ga'} \tag*{$(\ast)$}
\]

\uline{Mais c'est faux} \deleted{si} car on ne suppose pas $f$ et $g$ de degrés \inserted{pairs}

Les deux mb\unsure{rs} de la formule $(\ast)$ diffèrent par un facteur
\[
(-1)^{ag+bf+Bf}
\]

% L17–19 | lower middle
Si dans la définition de $k_{A\otimes B}$ on n'avait pas mis de signe $(-1)^{Bx}$, le signe qui \unsure{rendrait|viendraient} \unsure{correct|comme} $\unsure{(\ast)}$ serait
\[
(-1)^{a'g+bf}
\]

% L20–22 | lower
Si \inserted{de plus} on se dispensait dans $\unsure{(\ast)|\ast}$ de prendre $f_*\otimes g_*$ au sens \deleted{\unsure{gradué|produits}} tordu, et qu'on \deleted{\unsure{prenne|prend}} enlève le \uline{signe} $(-1)^{ga'}$, on trouverait donc un signe
\[
(-1)^{bf}=(-1)^{B'f}\unsure{.}
\]

% L23–24 | lower
Si on prenait $k$ tordu, $\otimes$ ordinaire, le signe serait
\[
(-1)^{bf+(AB+A'B')}
\]

% L25–27 | foot, boxed formula and note at right
C'est donc le signe où on prend tout ordinaire ($k_{A\otimes B}\inserted{=k_A\otimes k_B}$ et $f_*\otimes g_*$) qui est le plus simple. Donc on trouve alors
\[
\boxed{(f\times g)_*=(f_*\overset{\mathrm{ord}}{\otimes}g_*)\bigl(\mathrm{id}_A\otimes[(-1)^B\unsure{\eta|D}_{B'}]^{\deg f}\bigr)}
\]
\margin{\deleted{\illegible{1}} Il n'y a aucun signe dès que \textcircled{$f$} est pair.}
% archive sheet number: 28

\scanpage{30}
% sheet: reused_typed
% hand: typed
% L1 | top centre, in reading orientation
\mspage{-2.17-}

% L2–5 | top, end of the preceding proof
\[
\begin{aligned}
\lambda(\alpha)\cdot\gamma\cdot\lambda(\alpha^{-1})
 &=\handwritten{\chi}\bigl(\lambda(\alpha^{-1})\bigr)\cdot\gamma\\
 &=\alpha^{-\handwritten{\bar{\chi}}(\handwritten{\bar{\lambda}})}\cdot\gamma
\end{aligned}
\]
for $\gamma\in(\handwritten{\mathbb G}_a)_k$. Therefore $\lambda(\alpha)\cdot\gamma\cdot\lambda(\alpha^{-1})$ has a finite specialization when $\alpha\handwritten{\longrightarrow}0$, i.e. $\handwritten{\mathbb G}_a\handwritten{\subset}P(\lambda)$. Therefore $B\handwritten{\subset}P(\lambda)$. \uline{QED}

% L6–9 | upper middle
\uline{Proposition} \uline{2.7}: Let a reductive group $G$ act via $\sigma$ on a scheme $X$, proper over $k$. Then for all $x\in X_k$, $L\in\mathrm{Pic}^{G}(X)$, and 1-PS's $\lambda$ of $G$,
\[
\mu^{L}(x,\lambda)=\mu^{L}(x,\gamma^{-1}\cdot\lambda\cdot\gamma)
\]
\margin{\handwritten{\unsure{9}}}
if $\gamma\in P(\lambda)_k$.

% L10–14 | middle, start of proof and calculation
\uline{Proof}: Let $y$ be the specialization of $\sigma(\lambda(\alpha),x)$ as $\alpha\handwritten{\longrightarrow}0$. Suppose we write $y=\lim_{\handwritten{\alpha\to 0}}\sigma(\lambda(\alpha),x)$. In this notation, we calculate:
\[
\begin{aligned}
&\lim_{\handwritten{\alpha\to 0}}\sigma(\gamma^{-1}\cdot\lambda(\alpha)\cdot\gamma,x)\\
&\qquad=\lim_{\handwritten{\alpha\to 0}}\sigma\handwritten{\bigl(}\gamma^{-1}\cdot\lambda(\alpha)\cdot\gamma\cdot\lambda(\alpha^{-1}),\ \sigma(\lambda(\alpha),x)\handwritten{\bigr)}\\
&\qquad=\sigma\handwritten{\bigl(}\gamma^{-1}\cdot\lim_{\handwritten{\alpha\to 0}}\lambda(\alpha)\cdot\gamma\cdot\lambda(\alpha^{-1}),\ y\handwritten{\bigr)}
\end{aligned}
\]

% L15–19 | lower
Let $\lim_{\handwritten{\alpha\to 0}}\lambda(\alpha)\cdot\gamma\cdot\lambda(\alpha^{-1})=\gamma'$. According to Proposition 2.6, $\gamma'$ centralizes $\lambda$. Now $\mu(x,\gamma^{-1}\cdot\lambda\cdot\gamma)$ is calculated by looking at $\lim_{\handwritten{\alpha\to 0}}\sigma(\gamma^{-1}\cdot\lambda(\alpha)\cdot\gamma,x)$. Therefore:
\[
\mu(x,\gamma^{-1}\cdot\lambda\cdot\gamma)=\mu(\sigma(\gamma^{-1}\cdot\gamma',y),\gamma^{-1}\cdot\lambda\cdot\gamma)
\]

By property (i) of $\mu$ (\S1 above), we have:
% archive sheet number: 29

\scanpage{31}
% L1–2 | top
Pour la suite, on se borne aux $f=(\unsure{f^*|f'|f_*},f_*)$ \uline{de degré pair}, donc on trouve $(f\times g)_*\unsure{{:}|{=}}f_*\overset{\mathrm{ord}}{\otimes}g_*$.

% L3–8 | upper
Considérons l'isom de symétrie de $A\otimes B$ sur $B\otimes A$ [considérés comme produits cart. dans $\unsure{\mathcal V|V}$], on vérifie que c'est $s(x\otimes y)=(-1)^{xy}y\otimes x$. On a bien, comme il convient
\[
\begin{aligned}
(g\times f)_*\bigl(\underbrace{s(x\otimes y)}_{y\otimes x\,(-1)^{xy}}\bigr)
 &=s\bigl(\underbrace{(f\times g)_*(x\otimes y)}_{f_*(x)\otimes g_*(y)}\bigr)\\[-2pt]
\shortparallel &\\[-4pt]
g_*(y)\otimes f_*(x)\,(-1)^{xy}
 &\qquad g_*(y)\otimes f_*(x)\,(-1)^{(x+f)(y+g)}
\end{aligned}
\]

% L9–12 | centre, left square and pair at right
Prouvons les formules de projection
\diagram{\begin{tikzcd}[ampersand replacement=\&]
X\arrow[d,"f"'] \& X\times S\arrow[l,"q"']\arrow[d,"f'"]\\
Y\,\deleted{\unsure{\to}} \& Y\times S\arrow[l,"{p\deleted{\unsure{_*}}}"]
\end{tikzcd}}
\[
\left\{\begin{aligned}
p^*(f_*(x))&=f'_*(q^*(x))\\
f^*(p_*(u))&=q_*(f'{}^*(u))
\end{aligned}\right.
\]

% L13–17 | middle, lower square and workings on its right
\diagram{\begin{tikzcd}[ampersand replacement=\&,column sep=large]
x\in B\arrow[r,"{q^*(b)=b\otimes1}"] \& B\otimes\unsure{R|A}\\
A\arrow[u,"{f^*}"]\arrow[r,"{p^*(a)=a\otimes1}"'] \& A\otimes R\arrow[u,"{f^*\otimes\mathrm{id}_R=f'_*}"']
\end{tikzcd}}
\[
\begin{aligned}
 f_*(x)\otimes1&=(f_*\otimes1)(x\otimes1) &&\text{OK}\\
 f^*\bigl(\underbrace{(\mathrm{id}_A\otimes k_R)(a\otimes r)}_{\deleted{\unsure{a\,k(r)}}\,k(r)a}\bigr)
 &\overset{?}{=}(\mathrm{id}_B\otimes\unsure{k_R|\kappa_R})\underbrace{(f^*\otimes\mathrm{id}_R)(a\otimes r)}_{f^*(a)r}\\
 \shortparallel &\\[-3pt]
 k(r)f^*(a)&\qquad k(r)f^*(a)&&\text{O.K.}
\end{aligned}
\]

% L18 | lower, two chains above diagrams
\diagram{\begin{tikzcd}[ampersand replacement=\&]
X\arrow[r,"f"] \& Y\arrow[r,no head,"g"] \& Z \& A \& B\arrow[l,"{f^*}"'] \& C\arrow[l,"{g^*}"']
\end{tikzcd}}

% L19–22 | lower left
\diagram{\begin{tikzcd}[ampersand replacement=\&,column sep=small]
\&X\arrow[dl,"{\Gamma_{gf}}"']\arrow[dr,"{\Gamma_f}"]\&\\
X\times Z\arrow[dr,"{\Gamma_f\times\mathrm{id}_Z=\lambda}"']\& \&X\times Y\arrow[dl,"{\mathrm{id}_X\times\Gamma_g=\lambda'}"]\\
\&X\times Y\times Z\&
\end{tikzcd}}

% L19–22 | lower right, with element labels along arrows
\diagram{\begin{tikzcd}[ampersand replacement=\&,column sep=small]
\&\overset{\unsure{a}\,a f^*(b)\cdot f^*(g^*(c))}{A}\&\\
A\otimes C\arrow[ur] \& \& A\otimes B\arrow[ul]\\
\underset{\deleted{\unsure{a\otimes b}}}{\nwarrow}\&\underset{a\otimes b\otimes c}{A\otimes B\otimes C}\arrow[ul,"{\lambda^*}","{(a f^*(b))\otimes c}"']\arrow[ur,"{\lambda'{}^*}","{a\otimes b g^*(c)}"']\&
\end{tikzcd}}


% L23–24 | bottom left
\[
\lambda'{}^*(\lambda_*(u))=(\Gamma_f)_*(\Gamma_{gf}^*(u))
\]
et formule symétrique

% L23–27 | bottom right, set off by a rule
\margin{\uline{Question} Trouver un cas général commun où la formule d'échange soit valable. Suffit-il que le produit fibré reste dans $\mathcal V$ et que les dim. soient correctes ?}
% archive sheet number: 30

\scanpage{32}
% sheet: reused_typed
% hand: typed
% L1 | top centre, in reading orientation
\mspage{-2.13-}

% L2–7 | upper
\[
\phi=\psi_1\cdot\langle\lambda\rangle\cdot\psi_2
\]
for $R$-valued points $\psi_1$ and $\psi_2$ of $G$, and some 1-PS $\lambda$. Moreover $\lambda$ is non-trivial since $\phi$ is not an $R$-valued point of $G$. But then
\[
\xi'=\sigma(\psi_2,\xi)
\]
\[
\text{and}\qquad \sigma(\langle\lambda\rangle,\xi')=\sigma\handwritten{\bigl(}\psi_1^{-1},\sigma(\phi,\xi)\handwritten{\bigr)}
\]
are both $R$-valued points of $X$.
% L7–12 | upper middle, continuation
In other words, the $K$-valued, but not $R$-valued point $\langle\lambda\rangle\times\xi'$ of $G\times X$ is mapped by $\unsure{\Psi|\bar{\psi}}$ to an $R$-valued point of $X\times X$. Since by definition, $\langle\lambda\rangle$ is induced from a $K$-valued point of $\handwritten{\mathbb G}_m$ via $\lambda:\handwritten{\mathbb G}_m\handwritten{\longrightarrow}G$, this means that the composite morphism $\Psi\circ(\lambda\times1_X):\handwritten{\mathbb G}_m\times X\handwritten{\longrightarrow}X\times X$ is not proper, i.e. the induced action of $\handwritten{\mathbb G}_m$ is not proper. \uline{QED}

% L13–14 | centre
\uline{Corollary} \uline{2.5}: Let a reductive algebraic group $G$ act on an algebraic pre-scheme $X$. Let $L\in\mathrm{Pic}^{G}(X)$. Then $G$ acts properly on $X^s_{\handwritten{(\unsure{0|o})}}(L)$.

% L15–23 | lower, proof continuing off the sheet
\uline{Proof}: Without loss of generality, we may replace $X$ by $X^s_{\handwritten{(0)}}(L)$ : then $L$ is ample on $X$, and $X$ is a scheme. By the above proposition, it suffices to prove that for every non-trivial 1-PS $\lambda$ of $G$, the induced action of $\handwritten{\mathbb G}_m$ is proper. Fix some $\lambda$. Let $\handwritten{\unsure{\mathcal U|U}}$ be the open set in $X$ of points properly stable for the induced action of $\handwritten{\mathbb G}_m$: then I claim $\handwritten{\mathcal U}=X$. For if $x$ is a geometric point of $X$, then $x$ has an affine neighborhood $X_s$, where $s$ is an invariant section of $L^N$, and such that all stabilizers of points of $X_s$ are 0-dimensional. But then $s$ is invariant under $\handwritten{\mathbb G}_m$, and the action of $\handwritten{\mathbb G}_m$ on $X_s$
% archive sheet number: 31

\scanpage{34}
% L1–5 | top, title and definition
\uline{Catégorie des algèbres bivariantes} (sur un anneau $k$)

C'est la catégorie dont les objets sont les $k$-algèbres\inserted{\uline{\unsure{muni|munies|munie} d'une forme} $k_A:A\unsure{-|\to} k$}, et les morphismes de $A$ dans $B$ les couples $f=(f_*,f^*)$, où $f^*:B-A$ est un hom d'algèbres, et où $f_*:A-B$ est un hom $B$-linéaire\inserted{[\unsure{à|:} g. et \unsure{à|:} dr]}, i.e. un hom additif satisfaisant les ``form. de projection''

% L6–8 | upper, projection formulas and annotation
\[
\bigl[f_*(x f^*(y))=f_*(x)y\bigr],\quad f_*(\unsure{f|f^*}(y)x)=y f_*(x). \tag*{$(\ast)$}
\]
\margin{signe $(-1)^{\delta j}$ \unsure{raisonnable} ds le cas gradué, si la deux. formule est valable, $y$ de degré $j$, $f_*$ de degré $\delta$.}
% slanted
\margin{Il faut choisir entre l'une et l'autre formule !}
% slanted
et de plus :
\[
k_B f_*=k_A \tag*{$(\ast\ast)$}
\]

% L9–13 | upper middle
On a une variante graduée, \deleted{et} On aura à travailler dans diverses sous-catégories pleines, ou même non pleines.

Dans le cas gradué anticommutatif (où $f^*$ est donc un hom d'Alg. graduées, et \deleted{\unsure{de}} $f_*$ préserve à translation près \inserted{\uline{de degré pair}} les graduations), les deux formules $(\ast)$ sont équivalentes.

% L13 margin | left, slanted block
\margin{\uline{Remarque} \unsure{Pour|Bien} que $(f_*,f^*)$ soit un isomorphisme, il faut \unsure{et} que $f^*$ \deleted{soit un isom et} $f_*$ \unsure{bijectifs|bijectif}, i.e. $f^*$ un isom et $f_*(1)$ inversible ; les ``bons'' isom. sont ceux où $f_*(1)=1$, ou encore $f_*$ et $f^*$ sont inverses l'un de l'autre.}
% slanted

% L14–19 | middle
Soit $A$ une algèbre munie d'une forme
\[
k_A:A-k
\]
\deleted{que} \unsure{donc|ds|dans} la forme bilinéaire $\unsure{b|B}$
\[
(x,y)\longmapsto k(xy)
\]
sur $A\times A$ \deleted{est séparée, de} défini\unsure{t|\deleted{t}} un \inserted{(= deux !)} hom. \deleted{injectif}
\[
\varphi_A:A\longrightarrow\check A\qquad(\check A:\mathrm{Hom}_{k\text{-}\mathrm{mod}}(A,k)).
\]
\inserted{$\langle y,\varphi_A(x)\rangle=k(yx)$}
\margin{\deleted{\unsure{l'injection} \illegible{2}}}
% slanted

% L20–23 | lower middle
Supposons que $B$ soit de \deleted{\unsure{a}} muni de $k_B:B-k$, \inserted{tel que $\varphi_B$ \uline{soit injectif},} et soit $f^*:B\longrightarrow A$ tel que l'hom transposé $\,{}^{t}f^*:\check A-\check B$ applique $\mathrm{Im}\,\varphi_A$ dans $\mathrm{Im}\,\varphi_B$ [ce qui n'est pas une restriction si $\varphi_B:B\xrightarrow{\sim}\check B$ !)

Alors l'application induite, \inserted{notée} $f_*$, est donc caractérisée par
\[
k(y f_*(x))=k(f^*(y)x) \tag*{$(\ast\ast\ast)$}
\]

% L24–27 | foot, several additions along the lower edge
et satisfait de ce fait la $2^{\text{ème}}$ relation $(\ast)$, [donc la première dans le cas commutatif, ou gradué anticommutatif, quand $f_*$ est hom à translation \inserted{paire} près].
\inserted{De plus, il satisfait $(\ast\ast)$ (faire $y=1$ dans $(\ast\ast\ast)$).}
\deleted{\unsure{Supposons donc} $A,B$}
Réciproquement, $A$ étant \unsure{fixé}, \deleted{\unsure{étant}} tout hom $(f^*,f_*)$ de $(A,k_A)$ en $(B,k_B)$, $f_*$ satisfait $(\ast\ast\ast)$ ce qui prouve \unsure{que $k$ étant donné|$k$ étant donné}, $f_*$ est déterminé par $f^*$.
% archive sheet number: 33

\scanpage{35}
% sheet: reused_typed
% hand: typed
% L1 | top centre, in reading orientation
\mspage{-3.11-}

% L2–7 | upper
Moreover, the section:
\[
(D_{\handwritten{\alpha_0},\ldots\handwritten{\alpha_n}})^{n+1}\in\Gamma\handwritten{\bigl(}\handwritten{\unsure{\mathcal U|U}}_{\mathrm{reg}},\,L_{\handwritten{\alpha_0}}^{n+1}\handwritten{\otimes}\cdots\handwritten{\otimes}L_{\handwritten{\alpha_n}}^{n+1}\handwritten{\bigr)}
\]
is invariant in the $\mathrm{SL}(n+1)$-linearization of this sheaf, hence in the $\mathrm{PGL}(n)$-linearization of this sheaf, . Therefore, according to SGA 8, \S1, there is a section $E_{\handwritten{\alpha_0},\ldots\handwritten{\alpha_n}}$ of $M_{\handwritten{\alpha_0}}\handwritten{\otimes}\cdots\handwritten{\otimes}M_{\handwritten{\alpha_n}}$ such that:
\[
(D_{\handwritten{\alpha_0},\ldots\handwritten{\alpha_n}})^{n+1}=\Phi^*(E_{\handwritten{\alpha_0},\ldots\handwritten{\alpha_n}}).
\]

% L8–11 | upper middle
\fbox{\S2.} The purpose of this section is to show that although $M_0\handwritten{\otimes}\cdots\handwritten{\otimes} M_m$ is not ample on $Z$ - which is not even a scheme - it is ample on a fairly large open subset $Z_0\handwritten{\subset}Z$. We first define the subset $\Phi^{-1}(Z_0)$ of $\handwritten{\mathcal U}_{\mathrm{reg}}$ :

% L12–15 | middle
\uline{Definition} \uline{3.7}/\uline{Proposition} \uline{3.4}: $\handwritten{\mathcal U}_{\mathrm{stable}}$ is the open subset of $(\handwritten{\mathbb P}_n)^{m+1}$ whose geometric points $x=\bigl(x^{(0)},\ldots,x^{(m)}\bigr)$ are those points such that for every proper linear subspace $L\handwritten{\subset\overline{\mathbb P}}_n$ :
\[
\frac{\text{number of points }x^{(i)}\text{ in }L}{m+1}\handwritten{<}\frac{\dim L+1}{n+1}.
\]

% L16–21 | lower
\uline{Proof}: Consider the set of all pairs $\handwritten{\{}I,n_0\handwritten{\}}$ consisting of a subset $I\handwritten{\subset\{}0,1,\ldots,m\handwritten{\}}$ and an integer $n_0$ between $0$ and $n-1$ such that:
\[
\frac{\text{cardinality of }I}{m+1}\handwritten{\geq}\frac{n_0+1}{n+1}.
\]
Then the set of geometric points \uline{not} satisfying the above condition is the (finite) union over all $\handwritten{\{}I,n_0\handwritten{\}}$ of the sets consisting of
% archive sheet number: 34

\scanpage{36}
% L1–4 | top
graduées anticommutatives, $f^*$ \unsure{homogène|homogènes}. Supposons qu'il existe \deleted{deux} entier $d=d(A)$ tels que $k_A$ soit composé $A\unsure{-|\to}A^d\xrightarrow{\unsure{\bar k_A^{(d)}|k_A^{(d)}|\bar{\kappa}_A^{(d)}}} k$, et de même un $d'=d(B)$. Alors si $x$ est de degré $i$, la forme $\varphi_A(x)$ est nulle en degré $\ne d-i$,

% L5–8 | upper
donc $\langle f^*(y),\varphi_A(x)\rangle=k(f^*(y)x)$ est nul \deleted{\unsure{si}} en degré $\ne d-i$, \deleted{\unsure{donc si $f$|donc si $x$}} de $=$, \inserted{si $x'$ \deleted{\unsure{$=f_*(x)$}}, est hom de degré $i'$,} la forme $\varphi_B(x')$ est nulle en degré $\ne d'-i'$, donc on doit avoir (si $\unsure{f_*(x')|f_*(x)|\varphi_B(x')}\ne0$)
\[
d-i=d'-i'\quad\text{i.e.}\quad i'=i+(d'-d).
\]

% L8–12 | upper middle
Si alors on pose $f_*(x)=\sum_j x'_j$, on voit que l'effet de la forme $\varphi_{\unsure{B|A}}(f_*(x))$ sur $B^{d'-j}$ est égal à $\varphi_B(x'_j)$, donc est nul si $\deleted{j}\inserted{d'-j}\ne d-i$, donc dans ce cas $x'_j=0$. Donc

$f_*$ \uline{augmente les degrés de} $d'\unsure{-|+}d$.

% L13–18 | middle
Nous travaillons dans la catégorie \inserted{bivariante} des Algèbres graduées anticommutatives, munies d'un degré $d$, le $f_*:A\to B$ étant assujetti à être de degré $d(B)-d(A)$, \inserted{et $k_A$ \deleted{\unsure{de a|se}} factorisant par $A^{d(A)}$.} La sous catégorie pleine des $A$ tels que $d(A)$ pair est commode

% L19–24 | lower
[et il y a lieu parfois de se borner aux Algèbres $A$ qui sont à degrés \deleted{\unsure{entre}} dans $[0,d(A)]$]. Il n'est pas \deleted{\unsure{utile}} \inserted{commode} de se borner aux ``Algèbres de Poincaré'', associé\unsure{es} comme dessus à des formes sur $A$, car les constructions universelles qu'on va faire en feraient de \unsure{tte} façon sortir !
% archive sheet number: 35

\scanpage{37}
% sheet: reused_typed
% hand: typed
% L1 | top centre, in reading orientation
\mspage{-3.9-}

% L2–3 | top
leaves $x$ fixed. Secondly, $3)\handwritten{\Rightarrow}1)$ is an immediate consequence of Proposition 2.1.

% L4–11 | upper and middle
Thirdly, we will prove that $2)\handwritten{\Rightarrow}3)$. By virtue of 2), all the points $x_i$ cannot lie in one hyperplane, hence we can choose $n+1$ of the $x_i$ which are not in one hyperplane, say $x^{(\unsure{0|o})},x^{(1)},\ldots,x^{(n)}$. Without loss of generality, we may assume that $x^{(i)}$ has homogeneous coordinates $x_j^{(i)}=\delta_{ij}$. Now for each $n+k$ between $n+1$ and $m$, let $S_k'$ be the set of integers $i$ such that
\[
D_{0,1,\ldots,\handwritten{\hat}{i},\ldots,n,n+k}\handwritten{\ne}0,
\]
i.e. $x^{(n+k)}$ is not in the hyperplane spanned by $x^{(0)},\ldots,\handwritten{\widehat}{x^{(i)}},\ldots,x^{(n)}$.

% L12–18 | middle
Then I claim that there is no partition of the set $\handwritten{\{}0,1,\ldots,n\handwritten{\}}$ into 2 disjoint subsets $T'$ and $T''$ such that every $S_k'$ is contained in $T'$ or $T''$. For if there were, and if one let $L'$ (resp.$L''$) be the linear subspace defined by $X_i=0$ for all $i\in T'$ (resp. $i\in T''$), then every point $x^{(k)}$ would be in $L'\handwritten{\cup}L''$. It follows immediately that a suitable set of subsets $S_i\handwritten{\subset}S_i'$ is an $R$-partition $R$ and that $x\in\handwritten{\unsure{\mathcal U|U}}_R$. \uline{QED}

% L19–20 | lower
From this Proposition, it is clear that $\handwritten{\mathcal U}_{\mathrm{reg}}$ is precisely the set of pre-stable points of $(\handwritten{\mathbb P}_n)^{m+1}$, in the sense of Chapter 1.

% L21–25 | foot
Our next step is to construct a geometric quotient of $\handwritten{\mathcal U}_{\mathrm{reg}}$ by $\mathrm{PGL}(n)$: let $\handwritten{\mathcal U}_1,\ldots,\handwritten{\mathcal U}_N$ be the open subsets $\handwritten{\mathcal U}_R$ of $\handwritten{\mathcal U}_{\mathrm{reg}}$ and the subsets obtained from these by permuting the coordinates. Let $(Z_i,\phi_i)$ be the geometric quotient of $\handwritten{\mathcal U}_i$ by $\mathrm{PGL}(n)$. For all pairs $i,j$, $\handwritten{\mathcal U}_i\handwritten{\cap}\handwritten{\mathcal U}_j$ is an invariant open subset in $\handwritten{\mathcal U}_i$ and in $\handwritten{\mathcal U}_j$.
% archive sheet number: 36

\scanpage{38}
% L1–9 | top
\deleted{Soient $A$ et $B$ deux algèbres}

Le produit tensoriel gradué gauche $A\otimes B$ est un foncteur sur la catégorie \inserted{$\unsure{C|\mathcal C}$} des algèbres graduées\inserted{ \unsure{bivariante|bivariantes}}. Ce produit tensoriel (malheureusement !) ne s'interprète pas comme un produit dans la catégorie bivariante : car on n'a pas d'hom de la catégorie $f:A\to A\otimes A$, $f=(f^*,f_*)$, qui jouerait le rôle d'une application diagonale. On trouve néanmoins que c'est un produit \inserted{dans $C$} quand $A,B$ sont toutes deux des Alg. de Poincaré\unsure{\ldots}

% L10–15 | middle
Algèbre bivariante simpliciale : c'est un objet simplicial (\uline{pas} \uline{$1/2$} simplicial !) de la catégorie des Algèbres bivariantes, donc un foncteur \unsure{contravariant|contravar.} de la catégorie des ens finis non vides dans $C$. Il est dit \uline{spécial} ``s'il transforme sommes disjointes en produit tensoriel'', i.e. si

% L16–18 | lower middle
pour $I=I'\amalg I''$, les \deleted{\unsure{hom}}
$A(I')\xrightarrow{\alpha^*}A(I)$, $A(I'')\xrightarrow{\beta^*}A(I)$
[où $\alpha:I'\unsure{\to|-}I$, $\beta:I''-I$) \unsure{induisent|définissent} un \uline{isom}
\[
A(I')\otimes A(I'')\xrightarrow{\sim}A(I'\amalg I'').
\]

% L19–23 | lower
Si \deleted{$A(\Delta^0)$ est une algèbre} on se borne à la sous-catégorie pleine des Alg. biv. simpliciales $A(\unsure{*|\cdot})$ telles que $\unsure{A|\Delta}(\langle\unsure{\omega|0}\rangle)=A_0$ soit une algèbre de Poincaré, on trouve une \uline{équivalence} de la catégorie des Alg. biv. simpliciales et des alg. de Poincaré.
% archive sheet number: 37

\scanpage{39}
% sheet: reused_typed
% hand: typed
% L1 | top centre, in reading orientation
\mspage{-3.8-}

% L2–5 | upper
Our next aim is to show that we have enough open sets:

\uline{Definition} \uline{3.6}: For fixed $m$, let $\handwritten{\unsure{\mathcal U|U}}_{\mathrm{reg}}$ be the union of all the open subsets $\handwritten{\mathcal U}_R$ and their images under permutations of the factors of $(\handwritten{\mathbb P}_n)^{m+1}$.

% L6–11 | upper middle
\uline{Proposition} \uline{3.3}: Let $x=\handwritten{\bigl(}x^{(\unsure{0|o})},x^{(0)},\ldots,x^{(m)}\handwritten{\bigr)}$ be a geometric point of $(\handwritten{\mathbb P}_n)^{m+1}$. Then the following are equivalent:

(1) the stabilizer $S(x)$ is 0 dimensional,

(2) there do not exist disjoint proper linear subspaces $L'$ and $L''$ in $\handwritten{\overline{\mathbb P}}_n$ such that every $x^{(i)}$ is in either $L'$ or $L''$,

(3) $x$ is a geometric point of $\handwritten{\mathcal U}_{\mathrm{reg}}$.

% L12–18 | middle and lower
\uline{Proof}: Let $k$ be the algebraically closed field over which $x$ is defined. For simplicity, we shall write $\handwritten{\mathbb P}_n$ for $\handwritten{\mathbb P}_n\times k$, $\handwritten{\mathcal U}_R$ for $\handwritten{\mathcal U}_R\times k$, etc. in the course of this proof. First, the implication $1)\handwritten{\Rightarrow}2)$ is clear: for, in suitable homogeneous coordinates $\handwritten{\{}x_i\handwritten{\}}$, one may assume that:
\[
\begin{aligned}
L' &\handwritten{\subset\{}x_0=x_1=\ldots=x_r=0\handwritten{\}}\\
L'' &\handwritten{\subset\{}x_{r+1}=\ldots=x_n=0\handwritten{\}}.
\end{aligned}
\]

% L19–20 | foot, hand-drawn partitioned matrix
Then the subgroup of transformations:
\diagram{\begin{tikzcd}[ampersand replacement=\&,column sep=1.2em,execute at end picture={
\draw[dashed] (block.north) -- (block.south);
\draw[dashed] (block.west) -- (block.east);
\draw ($(block.north west)!0.16!(block.south east)$) -- ($(block.north west)!0.2!(block.south east)$);
\draw ($(block.north west)!0.26!(block.south east)$) -- ($(block.north west)!0.3!(block.south east)$);
\draw ($(block.north west)!0.61!(block.south east)$) -- ($(block.north west)!0.66!(block.south east)$);
\draw ($(block.north west)!0.72!(block.south east)$) -- ($(block.north west)!0.76!(block.south east)$);
\draw ($(block.north west)!0.8!(block.south east)$) -- ($(block.north west)!0.84!(block.south east)$);
\draw ($(block.north west)!0.87!(block.south east)$) -- ($(block.north west)!0.9!(block.south east)$);
}]
\handwritten{\begin{array}{c}\hline \uparrow\\[-4pt] r+1\\[-4pt]\downarrow\\ \hline \uparrow\\[-4pt] n-r\\[-4pt]\downarrow\\ \hline\end{array}}
\& |[alias=block,inner sep=0pt]| \handwritten{\left[\begin{array}{ccc@{\quad}ccc}
\alpha&&0&&&\\[-4pt]
&&&&0&\\[-4pt]
0&&\alpha&&&\\[3pt]
&&&\beta&&0\\[-4pt]
&0&&&&\\[-4pt]
&&\unsure{/}&&0&\beta
\end{array}\right]}
\end{tikzcd}}
% archive sheet number: 38

\scanpage{40}
% L1–2 | top
\uline{Proposition} Pour tout $n\geq1$, soit $A_n\subset H^{\unsure{*|\alpha}}(X^n)$ un ss-\deleted{\unsure{groupe}}\unsure{vectoriel|algèbre} \uline{stable} \inserted{par \uline{gpe sym.}}
Conditions équivalentes :

% L3–4 | upper
\deleted{\unsure{1}}(i) a) Pour tout \deleted{\unsure{$u\in H^*($}} $u\in A_{n+m}\subset H^*(X^n\times X^m)$ et $x\in A_n\in H^*(X^n)$, on a $u(x)\in H^*(X^m)$.

% L5–6 | upper, interlinear addition under c)
b) Si $x\in A_n$, $y\in A_m$, alors $x\otimes y\in A_{n+m}$ ; pour tt $n$, $1\in A_n$.

c) Pour \unsure{t\textsuperscript{e}} \deleted{\unsure{int}} surjection $[1,n+1]\unsure{\longrightarrow|\longleftrightarrow}[1,\unsure{n|2n}]$, son graphe ds $H^*(\unsure{X|\deleted{Y}X}^n\times X^{n+1})$ est $\in A_{2n+1}$, \inserted{Et $=$ chose pour tte injection $[1,n]\unsure{\leftarrow|-}[1,n+1]$}

% L7–9 | upper middle
\deleted{\unsure{2}}(ii) Les \deleted{\unsure{systèmes des}} $A_n$ sont des \uline{sous-algèbres}, stables par images inverses \inserted{et images directes} par \uline{projections} $X^{n+1}\longrightarrow X^n$, et \deleted{par images inverses} par images directes par injections $X^n\unsure{\leftarrow|\hookleftarrow}X^{n+1}$.

% L10–13 | middle
Ces conditions impliquent que les $A_n$ sont stables pour l'opér. images inv. et images dir. par tte application simpliciale, \inserted{$X^n-X^m$} et que \deleted{$A_n$} pour tte telle application, $H^*(X^n\times X^m)$ contient la classe du graphe....

% L14–18 | lower written area
(iii) Considérons $A_{n+m}\in H^*(X^n\times X^m)$ comme un ens d'hom de $H^*(X^n)$ en $H^*(X^m)$, on veut que ces ens. soient stables par composition $(n,m\geq0)$, contiennent les identités, soient stables par $\unsure{\boxtimes|\otimes|\blacksquare}$, et contiennent les \unsure{homs} images inverses et directes par une application simpliciale.
% archive sheet number: 39

\scanpage{41}
% sheet: reused_typed
% hand: typed
% L1 | top centre, in reading orientation
\mspage{-3.13-}

% L2–3 | upper
The key property of stable geometric points turns out to be the following:

% L4–9 | upper and middle
\uline{Proposition} \uline{3.6}: Let $x$ be a geometric point of $(\handwritten{\mathbb P}_n)^{m+1}$. Then $x$ is stable if andonly if there are integers $N\geq N_0>0$ and there are monomials $P_0,\ldots,P_m$ in the expressions $D_{\handwritten{\alpha_0},\ldots,\handwritten{\alpha_n}}$ such that, for $0\leq i\leq m$:

(i) $P_i$ is a section of $L_0^N\handwritten{\otimes}\cdots\handwritten{\otimes}L_i^{N-N_0}\handwritten{\otimes}\cdots\handwritten{\otimes}L_m^N$

(ii) $P_i(x)\ne0$

% L10–14 | middle
\uline{Proof}: Let $E=(\handwritten{\mathbb R})^{m+1}$ be Euclidean space of dimension $m+1$. To each $D_{\handwritten{\alpha_0},\ldots,\handwritten{\alpha_n}}$ associate a point $d_{\handwritten{\alpha_0},\ldots\handwritten{\alpha_n}}\in E$ as follows:
\[
\begin{aligned}
d_{\handwritten{\alpha_0},\ldots\handwritten{\alpha_n}}&=(\unsure{\epsilon|\varepsilon}_0,\epsilon_1,\ldots,\epsilon_m)\\
\epsilon_i&=0\quad\text{if }\ i\handwritten{\notin\{}\handwritten{\alpha_0},\ldots\handwritten{\alpha_n}\handwritten{\}}\\
\epsilon_i&=1\quad\text{if }\ i\handwritten{\in\{}\handwritten{\alpha_0},\ldots\handwritten{\alpha_n}\handwritten{\}}
\end{aligned}
\]

% L15–20 | lower
Then $d$ determines which invertible sheaf $D$ is a section of. More generally, to any monomial $P$ is the $D$'s, associate the corresponding sum of the $d$'s. Now define $C_x$ to be the convex cone spanned by those points $d_{\handwritten{\alpha_0},\ldots,\handwritten{\alpha_n}}$ such that $D_{\handwritten{\alpha_0},\ldots,\handwritten{\alpha_n}}(x)\ne0$, i.e. such that the components $x^{(\handwritten{\alpha_0})},\ldots,x^{(\handwritten{\alpha_n})}$ do not lie in a hyperplane. Then it is easy to see that:
% archive sheet number: 40

\scanpage{42}
% L1–2 | top, small note above the bracket
Catégorie des $k$-Algèbres $[$\uline{graduées}$]$
% slanted
\margin{\unsure{sera souvent|pas néc.} \unsure{sous-entendu|ss.-entendu}} bivariantes (\inserted{anticommutatives,} munies d'un degré $d(A)$, nulles en degré $\notin[1,d(A)]$ ; on ne donne pas de $k_A:A-k$).

% L3–5 | upper
Objet particulier de cette catégorie : $k$ de degré 0.

Algèbre graduée bivariante jaugée : munie d'un hom dans $k$. Dans certains cas d'hom de telles, $f_*$ est déterminé par $f^*$. $(*)$

% L6–8 | upper
Produit tensoriel d'Algèbres bivariantes (resp. jaugées). Dans certains cas, c'est un produit dans la catégorie des Algèbres bivariantes jaugées.

% L9–12 | upper middle, between and beneath lines
Algèbre graduée bivariante \inserted{(resp. gr. biv. jaugée)} \uline{simpliciale}. \deleted{\unsure{Exemp}} Cas \uline{spécial} : on a $A(I'\amalg I'')\xleftarrow[\sim]{(\alpha^*,\beta^*)}A(I')\otimes A(I'')$, d'autre part si $\alpha:I-J$, $\beta:I'\to J'$, alors $\alpha\amalg\beta=\unsure{\gamma|r}:I\amalg I'\longrightarrow J\amalg J'$, on veut que $(\alpha\amalg\beta)_*:A(\unsure{J\amalg J'|I\amalg I'})-A(I\amalg I')$ s'identifie à $\alpha_*\otimes\beta_*$.

% L12–17 | middle
On voit que si $A(\Delta^0)=A_0$ est ``de Poincaré'' (cas jaugé) $A(\deleted{\unsure{*}})-A(0)$ est une équiv. des $A(*)$ \unsure{sur spéc|sur spés} de Poincaré, et des alg. de Poincaré. Et $A(*)-(A(0),\delta)$ est un foncteur pl. fid de la catégorie des Alg. graduées bivariantes jaugées simpliciales, dans la catégorie des algèbres \inserted{gr.} bivariantes jaugées $A$, munies d'un $\Delta\in[A\otimes A]^{d(A)}$ : $[$il convient de préciser les axiomes à mettre sur ce $\Delta]$
% slanted
\margin{\uline{à vérifier} ; il faut d'abord noter qu'il suffit de connaître $\delta_*:A\longrightarrow A\otimes A$ ; on prend $\Delta=\delta_*(1)$}

% L18–25 | lower, example and left margin
\uuline{Ex} 1) Si on a une théorie ch. \deleted{sur la \unsure{variété}} $[$satisfaisant Künneth et Poincaré$]$ sur la catégorie $\unsure{\mathcal V|V}$\margin{$k$ un \uline{corps}} des variétés complètes \unsure{équidim.} non singulières de dim $d$\inserted{, sur un corps donné}\deleted{, \unsure{de dim}, \illegible{1}} d'où un foncteur \inserted{cov} de $\mathcal V$ dans la catégorie $C_{\mathrm{jaug}}$ des Alg. gr. bivar. jaugées sur $k$, et en transformant par ce foncteur l'objet simplicial canonique $I\longmapsto X^I$ défini par $X\in\mathcal V$, on trouve une Algèbre \deleted{bivariante} simpliciale \deleted{bivariante} jaugée canonique, d'ailleurs définie aussi canoniquement, comme dit plus haut, par \unsure{*|le} \unsure{uniquement|\deleted{correspondants}|composant} $H^*(X)$\ldots.

% L26–27 | footnote under horizontal rule
$(*)$ Attention, les \uline{isom} $(f^*,f_*)$ ne sont pas néc. tels que $f_*=f^{*-1}$ i.e. $f_*(1)=1$ ($\Rightarrow$ \deleted{\unsure{tels}} que $k_A f^*=k_B$), ex : $f^*=\mathrm{id}$, $f_*(y)=y\unsure{v|\upsilon|\sigma}$ $(v\in A^*)$, $k_B(y)=k_A(yv^{-1})$ \ldots
% archive sheet number: 41

